
Permite desplazamiento horizontal y giro. Impide desplazamiento vertical.

\safesection{Articulación}
En dos dimensiones:
\[
\text{2 reacciones},\qquad R_x,\ R_y.
\]
Impide las dos traslaciones y permite el giro.

\safesection{Empotramiento}
En dos dimensiones:
\[
\text{3 reacciones},\qquad R_x,\ R_y,\ M.
\]
Impide
\[
u_x=0,\qquad u_y=0,\qquad\theta=0.
\]

\safesection{Esquema}
\begin{figure}[htbp]
\centering
\begin{tikzpicture}[x=0.90cm,y=0.90cm,>=Latex,
  force/.style={-{Latex[length=2.3mm]},BookBlue,line width=1.0pt},
  support/.style={draw=BookGray,line width=.95pt},
  beamline/.style={draw=BookGray,line width=2.0pt}
]
% -------------------- APOYO MÓVIL --------------------
\begin{scope}[xshift=0cm]
  \draw[beamline] (0.2,2.55)--(3.2,2.55);
  \draw[support] (1.00,2.55)--(0.62,2.03)--(1.38,2.03)--cycle;
  \draw[support] (0.72,1.86) circle (0.10);
  \draw[support] (1.28,1.86) circle (0.10);
  \draw[support] (0.40,1.66)--(1.60,1.66);
  \foreach \x in {0.50,0.80,1.10,1.40}{
    \draw[BookGray!60,line width=.55pt] (\x,1.66)--(\x-0.22,1.45);
  }
  \draw[force] (1.00,0.72)--(1.00,1.72);
  \node[font=\small,text=BookBlue,right] at (1.05,1.10) {$R_y$};
  \node[font=\sffamily\small\bfseries,text=BookBlue] at (1.7,3.18) {Apoyo móvil};
  \node[font=\sffamily\scriptsize,text=BookGray,align=center] at (1.7,0.15)
    {impide una\\traslación};
\end{scope}

% -------------------- ARTICULACIÓN --------------------
\begin{scope}[xshift=5cm]
  \draw[beamline] (0.2,2.55)--(3.2,2.55);
  \draw[support] (1.00,2.55)--(0.52,1.88)--(1.48,1.88)--cycle;
  \draw[support] (0.35,1.68)--(1.65,1.68);
  \foreach \x in {0.50,0.80,1.10,1.40}{
    \draw[BookGray!60,line width=.55pt] (\x,1.68)--(\x-0.22,1.47);
  }
  \draw[force] (1.00,0.72)--(1.00,1.78);
  \draw[force] (-0.05,2.55)--(0.86,2.55);
  \node[font=\small,text=BookBlue,right] at (1.05,1.10) {$R_y$};
  \node[font=\small,text=BookBlue,above] at (0.35,2.57) {$R_x$};
  \node[font=\sffamily\small\bfseries,text=BookBlue] at (1.7,3.18) {Articulación};
  \node[font=\sffamily\scriptsize,text=BookGray,align=center] at (1.7,0.15)
    {impide dos\\traslaciones};
\end{scope}

% -------------------- EMPOTRAMIENTO --------------------
\begin{scope}[xshift=10cm]
  \draw[beamline] (0.35,2.55)--(3.35,2.55);
  \draw[BookGray,line width=1.8pt] (0.35,1.35)--(0.35,3.45);
  \foreach \y in {1.50,1.80,...,3.30}{
    \draw[BookGray!60,line width=.55pt] (-0.02,\y-0.22)--(0.35,\y);
  }
  \draw[force] (0.35,0.72)--(0.35,1.95);
  \draw[force] (-0.90,2.55)--(0.22,2.55);
  \draw[BookBlue,line width=1.0pt,-{Latex[length=2.3mm]}]
    (1.05,3.15) arc[start angle=55,end angle=-215,radius=0.56];
  \node[font=\small,text=BookBlue,right] at (0.40,1.10) {$R_y$};
  \node[font=\small,text=BookBlue,above] at (-0.40,2.58) {$R_x$};
  \node[font=\small,text=BookBlue] at (1.35,3.05) {$M$};
  \node[font=\sffamily\small\bfseries,text=BookBlue] at (1.85,3.82) {Empotramiento};
  \node[font=\sffamily\scriptsize,text=BookGray,align=center] at (1.85,0.15)
    {impide traslaciones\\y giro};
\end{scope}
\end{tikzpicture}
\caption{Vinculaciones exteriores habituales en un problema plano y reacciones asociadas.}
\label{fig:vinculaciones-planas}
\end{figure}

\begin{table}[htbp]
\centering
\caption{Grados de libertad restringidos y reacciones de las vinculaciones planas básicas.}
\label{tab:vinculaciones-planas}
\begin{tabular}{lccc}
\toprule
Vinculación & Traslación $x$ & Traslación $y$ & Giro $\theta$\\
\midrule
Apoyo móvil & libre & impedida & libre\\
Articulación & impedida & impedida & libre\\
Empotramiento & impedida & impedida & impedido\\
\bottomrule
\end{tabular}
\end{table}

\begin{tip}
Antes de escribir una sola ecuación de equilibrio, \textbf{sustituye todos los apoyos por sus reacciones}. Ese dibujo es tu verdadero problema mecánico.
\end{tip}

% ============================================================
\safechapter{Esfuerzos internos}
% ============================================================

\safesection{Concepto}
Realizamos un corte en una barra.

Para que cada una de las dos partes continúe en equilibrio debemos introducir las acciones que la parte eliminada ejercía sobre la parte conservada.

Estas resultantes se denominan \textbf{esfuerzos internos}.

En tres dimensiones pueden aparecer:
\[
N_x,\qquad C_y,\ C_z,\qquad M_x,\ M_y,\ M_z.
\]

